% !TeX root = ../../../main.tex
\begin{ChThreeAbstract}
This study quantifies the association between non-adherence to antipsychotic medications and adverse outcomes in individuals with schizophrenia. 
We frame the problem using survival analysis, focusing on the time to the earliest of several adverse events (early death, involuntary hospitalization, jail booking).
We extend standard causal inference methods (T-learner, S-learner, nearest neighbor matching) to utilize various survival models to estimate individual and average treatment effects, where treatment corresponds to medication non-adherence. 
Analyses are repeated using different amounts of longitudinal information (3, 6, 9, and 12 months). 
Using data from Allegheny County in western Pennsylvania, we find strong evidence that non-adherence advances adverse outcomes by approximately~1 to~4 months. 
Ablation studies confirm that county-provided risk scores adjust for key confounders, as their removal amplifies the estimated effects. 
Subgroup analyses by medication formulation (injectable vs.~oral) and medication type consistently show that non-adherence is associated with earlier adverse events. 
These findings highlight the clinical importance of adherence in delaying psychiatric crises and show that integrating survival analysis with causal inference tools can yield policy-relevant insights. 
We caution that although we apply causal inference, we only make associative claims and discuss assumptions needed for causal interpretation.
\end{ChThreeAbstract}



% \vspace{-1em}
\ChThreeParagraph{Data and Code Availability}
This paper utilizes administrative data from the Allegheny County Department of Human Services in the United States. 
% \textcolor{red}{For the purposes of peer review, we have anonymized the specific county and the full department name.} The dataset consists of de-identified records, where part of the de-identification process involved coarsening the time resolution to the month level. %, %, where personally identifiable information (PII) has been removed,
% and its time resolution has been intentionally coarsened to the month level.
Due to confidentiality agreements, the dataset is not publicly available.

The implementation of our proposed method is available at our GitHub repository: \linebreak
\url{https://github.com/Shahriarnz14/causal-meta-learner-survival-analysis}.

% \vspace{-2em}
\ChThreeParagraph{Institutional Review Board (IRB)}
This study does not require IRB approval, as it conducts secondary analysis of de-identified data provided by the Allegheny County Department of Human Services.






% \vspace{-1em}
